\documentclass[12pt,a4paper,twocolumn]{article}
\usepackage[utf8]{inputenc}
\usepackage[spanish]{babel}
\usepackage{amsmath}
\usepackage{amssymb}
\usepackage{physics}
\usepackage{geometry}
\geometry{margin=0.5cm}

\title{Formulario de Ecuaciones Teoría Electromagnética Tema 2}
\author{Emerson Warhman}
\date{}

\begin{document}
\small
\maketitle
\section*{Ley de Biot-Savart}

La fuente de un campo magnético estable puede ser un imán permanente, un campo eléctrico variable o una corriente eléctrica directa.

\begin{equation*}
    dH = \frac{I \, dl \times \hat{r}}{4\pi r^2} [A/m]
\end{equation*}

\begin{equation*}
    H = \int dH = \frac{I}{4\pi} \int \frac{dl \times \hat{r}}{r^2} [A/m]
\end{equation*}

Las formas alternas de la ley de Biot-Savart son:
\begin{equation*}
    H=\int_s \frac{K \times \hat{n}}{r^2} dA [A/m]
\end{equation*}

\begin{equation*}
    H=\int_v \frac{J \times \hat{r}}{r^2} dv [A/m]
\end{equation*}

donde J es la densidad de corriente y K es la densidad de corriente superficial.

\end{document}
